\section{A nonreset, noncommuting exact instrument realization}\label{sec:supp-interior}
Here is an explicit small realization of the branch-interiority theorem. Let the core vertices be the six Pauli eigenstates
\[
 \tau_{j,\pm}=(I\pm\sigma_j)/2,\qquad j=x,y,z.
\]
Their convex hull is the Bloch octahedron. For any positive matrix $B$ with trace $t$ and unnormalized Bloch coordinates $b_j=\tr(\sigma_jB)$, if $\sum_j|b_j|\le t$, then
\begin{equation}\label{eq:octahedral-weights}
 B=\sum_{j,\pm}w_{j,\pm}\tau_{j,\pm},\qquad
 w_{j,\pm}=\frac{t-\sum_l|b_l|}{6}+\max\{\pm b_j,0\}.
\end{equation}
All coefficients are nonnegative, their sum is $t$, and their signed coordinate differences are $b_j$, proving the identity directly. If $B=0$ every coefficient is zero. Thus this formula uses no undefined zero-report posterior.

Take any qubit instrument and, after each reported branch, physically apply
$\mathcal D(B)=B/4+3\tr(B)I/8$. Its normalized output has Bloch length at most $1/4$; by Cauchy--Schwarz its Bloch coordinate sum is at most $\sqrt3/4<1$. Formula~\eqref{eq:octahedral-weights} gives exact branch rows with six core states for any finite command family. A separately specified initial state can be retained as one extra anchor. This seven-state upper is better for this example than the general dimension-only bound; it is not asserted minimal.

The instruments used in the exact diagnostics include amplitude damping with Kraus matrices
\[
 A_0=\operatorname{diag}(1,3/5),\qquad
 A_1=\begin{pmatrix}0&4/5\\0&0\end{pmatrix},
\]
and its conjugate by the Hadamard unitary. Their report probabilities depend on the input state and may vanish. After the declared depolarization their unconditional channels have distinct fixed Bloch vectors $(0,0,16/91)$ and $(16/91,0,0)$, so there is no common stationary density matrix. Their report branches are not all replacement maps: in particular the invertible branch retains a nonzero dependence on the input. The main theorem therefore has content beyond a shared fixed point or a reset model.

Two additional weak-measurement commands, with Kraus matrices
$\operatorname{diag}(3/5,4/5)$ and $\operatorname{diag}(4/5,3/5)$ and their Hadamard conjugates, have both report effects bounded away from zero. A two-label observer alternates these commands after report zero and stops at report one. Its physical block cost and score potentials are computed independently of the seven-state simulator's fourteen-state joint killed matrix. Exact rational inversion gives the same initial cost and squared-audit risk. This is a diagnostic of the all-policy interpolation proof, not its substitute.

\section{Interiority can first occur at a genuine second call}\label{sec:supp-block-interior}
Let one command have the single reported channel $\Phi(B)=(HBH^*+SBS^*)/2$, where $H$ is Hadamard and $S=\operatorname{diag}(1,i)$. The Kraus choice is not a recorded report. On Bloch vectors its matrix is
\[
 A=\begin{pmatrix}0&-1/2&1/2\\1/2&-1/2&0\\1/2&0&1/2\end{pmatrix}.
\]
The largest singular value of $A$ is one, so no positive uniform one-call floor holds. But $\|A^2\|_{\rm F}^2=5/8<49/64$, hence every output of $\Phi^2$ has Bloch length at most $7/8$ and spectral floor at least $I/16$. Thus the block theorem applies with $L=2$ and $\eta=1/8$. The realization still exposes both actual calls and allows stopping after the first. Treating the two unrecorded Kraus choices as if they were available reports would be a different instrument.

The average-interiority counterexample in the main article goes in the opposite direction: the random unitary choice is actually recorded, so its branch states remain pure and its record cannot be averaged away. These two examples fix exactly which information is erased and which remains in the causal interface.

\section{Positive coordinates, arbitrary effects and a forbidden inference}\label{sec:supp-positive-scope}
A state-positive simulator has a linear reconstruction map $p\mapsto\sum_ip_i\tau_i$ from its simplex. It does not have a positive affine inverse defined on every quantum state. If such an inverse $E:\mathcal D_d\to\Delta_M$ reproduced every density matrix exactly, then every pure state would be a convex combination of the fixed $\tau_i$. Extremality would force that pure state to be one of the finitely many $\tau_i$, an impossibility for $d\ge2$. Our construction instead fixes finitely many initial preparations and uses the strict branch margin after the first call. It is not an entanglement-preserving channel simulation.

Theorems~\ref{thm:r68-instrument} and~\ref{thm:r68-noise-exponent} use positive density coordinates. The lower count for uniform noisy states therefore cannot be imported into the general optimal-face realization problem without checking that its generator decoders have this structure. Conversely the fixed-angle spectral exclusion in Corollary~\ref{cor:r68-qubit-lower} requires only a finite stochastic matrix and a probability decoder, so it does not have that extra restriction. Its asymptotic lower rate is much weaker than the state-positive noise exponent and is not claimed sharp.

Fanizza, Lumbreras and Winter's Theorem 2 \cite{Fanizza24} bounds the maximal positive-realization dimension across a varying phase family in a different noisy quantum process. It includes an inverse-square-root noise dependence in that worst-phase statement. The present fixed-angle rational exclusion should not be presented as an improvement to it. The new uniform upper instead supplies a constructive finite realization from a condition on every actual report branch, and its same-policy corollary keeps the simulator hidden under control and stopping.
